\chapter{Estado del arte}
\label{chap:estado-arte}

Este capítulo recorre la literatura que el trabajo necesita, en el orden en que la va a usar: las
regularidades factoriales que orientan el diseño de los agentes y sirven de baseline, lo que el
aprendizaje automático añade y lo que arriesga, los sesgos que condicionan cualquier evidencia de
esta clase, y por qué la métrica de selección es el Rank-IC y no la rentabilidad. Cierra declarando
cómo se posiciona el TFM ante las tensiones que esa literatura deja abiertas.

\section{De las anomalías a las variables explicativas}

La literatura de factores parte de una observación sencilla: los retornos esperados no parecen
independientes de todas las características observables de una empresa. El punto de partida
histórico suele situarse en Banz (1981), que documentó que las empresas de menor capitalización
obtenían retornos ajustados por riesgo superiores a los que predecía el modelo de mercado: una
característica contable y pública, disponible para cualquiera, quedaba asociada a una diferencia
sistemática de rentabilidad. Los modelos de Fama y French formalizaron después primas asociadas a
tamaño y valor (Fama y French, 1992, 1993); trabajos
posteriores mostraron persistencia del momentum (Jegadeesh y Titman, 1993) y una prima ligada
a calidad (Asness et al., 2019). Carhart (1997) incorporó el momentum como cuarto factor
al modelo de tres, y su trabajo es relevante aquí por una razón adicional: mostró que buena parte
de la persistencia aparente en el rendimiento de los fondos de inversión desaparecía al controlar
por exposiciones factoriales conocidas. Es exactamente la lógica que este TFM aplica a su propia
señal en el análisis de neutralización de estilo del Capítulo~\ref{chap:resultados}.

Estas relaciones no deben entenderse como leyes mecánicas.
Su magnitud cambia con el horizonte, el universo, los costes y el régimen, y una señal que funciona
en promedio puede atravesar años adversos. El propio efecto tamaño de Banz se debilitó de forma
notable en las décadas posteriores a su publicación, un recordatorio de que una regularidad
documentada no es una garantía de rendimiento futuro.

La ampliación posterior del catálogo factorial hace más importante esa cautela. Novy-Marx (2013)
mostró que la rentabilidad bruta contiene información comparable a valor; Fama y French (2015)
incorporaron rentabilidad e inversión a su modelo de cinco factores; y Hou et al. (2015) organizaron
varias anomalías mediante una lectura basada en inversión. Frazzini y Pedersen (2014), por su parte,
relacionaron la baja beta con restricciones al apalancamiento. Estas extensiones no producen una
lista indiscutible de primas: producen familias de explicaciones que comparten variables, cambian de
magnitud con la muestra y pueden describir una misma acción de formas distintas.

La literatura de replicación convirtió esa proliferación en una pregunta empírica. McLean y Pontiff
(2016) compararon 97 predictores antes y después de su publicación y documentaron una caída apreciable
fuera de muestra y otra adicional cuando el mercado pudo explotar el hallazgo. Hou et al. (2020)
reevaluaron 452 anomalías bajo procedimientos homogéneos y mostraron que muchas dejan de superar
umbrales más exigentes o dependen de microcapitalizaciones. La implicación para este trabajo no es
que toda regularidad sea falsa, sino que una señal aprendida debe demostrar estabilidad temporal,
resistir controles conocidos y declarar qué parte de su aparente éxito procede de haber buscado entre
alternativas.

Para este TFM los factores cumplen tres papeles. Primero, orientan la división de features entre
agentes: rentabilidad y solvencia para \textit{quality}; ratios de valoración para \textit{value};
aceleración de resultados para \textit{growth}; tendencias para \textit{momentum}; y medidas de
volatilidad y liquidez para \textit{risk}. Segundo, sirven de baseline: una fórmula determinista de
un factor debe compararse contra el sistema aprendido sobre el mismo panel. Tercero, son controles
de atribución: si la señal desaparece al neutralizar exposiciones conocidas, la interpretación debe
ser más modesta.

No se afirma que un agente represente un factor puro. Por ejemplo, una empresa de alta calidad puede
presentar al mismo tiempo baja volatilidad, crecimiento persistente y una valoración elevada. La
especialización organiza la entrada del modelo; no impone que el mercado recompense cada bloque por
separado ni evita interacciones posteriores dentro del LightGBM.

\section{Aprendizaje automático en \textit{asset pricing}}

Los métodos de aprendizaje automático amplían la capacidad de modelar no linealidades e
interacciones entre características. Gu et al. (2020) muestran que estos modelos pueden
mejorar la predicción transversal cuando el conjunto de señales es amplio y el entrenamiento respeta
la estructura temporal. LightGBM, usado aquí en cada agente, pertenece a la familia de árboles
potenciados: construye de forma secuencial árboles débiles que corrigen errores de los anteriores y
puede manejar relaciones no monotónicas (Ke et al., 2017).

Ese resultado forma parte de una literatura más amplia que intenta reducir dimensionalidad sin
suponer que cada característica actúa linealmente y de forma aislada. Kelly et al. (2019) vinculan
características y covarianzas mediante componentes principales instrumentados; Freyberger et al.
(2020) combinan selección regularizada y funciones no paramétricas; y Feng et al. (2020) estudian
qué factores aportan información incremental dentro de un catálogo amplio. Sus métodos son
distintos, pero convergen en dos ideas relevantes aquí: la no linealidad puede importar y añadir
predictores no garantiza añadir información.

La elección de árboles potenciados no se presenta por ello como una conclusión universal. Frente a
un modelo lineal, ganan capacidad para representar umbrales e interacciones, pero pierden una
interpretación directa de coeficientes. Frente a una red profunda, ofrecen una complejidad más
contenida y una atribución local más manejable para el tamaño del panel. El catálogo conserva una
alternativa lineal regularizada y una formulación de ranking directo precisamente para que estas
preferencias puedan contrastarse en vez de asumirse.

Hay además una diferencia entre predecir retornos y ordenar acciones. Un error pequeño en unidades
de retorno puede convivir con un orden transversal pobre, y una transformación monótona del score
puede alterar el error cuadrático sin cambiar ninguna decisión de selección. Este trabajo sigue la
segunda pregunta: cada agente aprende una representación continua, pero el protocolo juzga el orden
que induce. Esa elección acerca el objetivo estadístico al uso final sin hacer que la cartera entre
prematuramente en la selección.

Esa flexibilidad tiene un coste epistemológico. Un modelo flexible puede encontrar patrones espurios
con facilidad, especialmente si se prueban parámetros hasta maximizar una métrica final. Por ello la
pregunta no es si un algoritmo complejo puede ajustar la historia, sino si conserva una ventaja al
evaluarse en observaciones posteriores no utilizadas para su entrenamiento ni para su selección.
El diseño \textit{walk-forward}, las etiquetas cerradas y la era reservada son respuestas a esa
exigencia. La Tabla~\ref{tab:posicionamiento} enfrenta las cuatro líneas de literatura que este
trabajo cruza con el riesgo metodológico de cada una y con lo que aquí se hace al respecto.

\begin{table}[H]
\centering
\small
\caption{Posicionamiento metodológico del estudio frente a líneas de literatura relevantes.}
\label{tab:posicionamiento}
\begin{tabularx}{\textwidth}{lXXX}
\toprule
Enfoque & Aportación principal & Riesgo metodológico & Respuesta de este TFM\\
\midrule
Factores clásicos & Interpretabilidad económica y baselines simples & Confundir una prima histórica con señal universal & Comparación sobre el mismo panel y neutralización de estilo\\
ML transversal & Interacciones y no linealidades & Sobreajuste de parámetros y dependencia temporal & \textit{Walk-forward}, catálogo cerrado y Rank-IC OOS\\
Backtest de cartera & Traducción a utilidad económica & Elegir por CAGR, ignorar costes o supervivencia & Cartera posterior al ganador y costes explícitos\\
Inferencia con múltiples pruebas & Penalización del \textit{data snooping} & Publicar sólo la mejor especificación & Bootstrap, permutación, placebos y Deflated Sharpe\\
\bottomrule
\end{tabularx}
\caption*{Elaboración propia a partir de Gu et al. (2020), Harvey et al. (2016) y
López de Prado et al. (2014).}
\end{table}

\section{Sesgos que condicionan la evidencia}

El \textit{lookahead bias} aparece cuando una variable incorpora información que no era conocida en
la fecha de decisión. En fundamentales, usar la fecha de cierre fiscal en lugar de la fecha de filing
es un ejemplo clásico: el trimestre puede haber terminado, pero sus resultados no se publicaron
aún. El sesgo de supervivencia surge cuando se usa la composición actual de un índice para describir
el pasado; se eliminan así empresas que salieron del universo y se exagera la calidad de la muestra
(Brown et al., 1992).

El \textit{data snooping} es distinto: incluso con datos perfectamente cerrados, muchas pruebas
incrementan la probabilidad de encontrar una aparente ganadora. Harvey et al. (2016) alertan
del gran número de factores propuestos en la literatura y de la necesidad de elevar el estándar de
evidencia. El Deflated Sharpe Ratio (Bailey y López de Prado, 2014) es útil porque descuenta la
selección entre configuraciones, pero
no sustituye a la disciplina de congelar una era reservada ni a contrastes de permutación y placebo.

White (2000) y Hansen (2005) formalizan dos respuestas relacionadas al problema de escoger el mejor
modelo después de comparar muchos: el \textit{Reality Check} y el test de capacidad predictiva
superior. Bailey et al. (2017) expresan la misma preocupación como probabilidad de sobreajuste del
backtest. Ninguna corrección reconstruye una hipótesis verdaderamente previa después de observar los
datos; por eso el catálogo cerrado, el ledger de decisiones y la reserva cumplen una función que no
puede reemplazarse con un único $p$-valor.

El sesgo de supervivencia tampoco se reduce a incluir o excluir empresas quebradas. Shumway (1997)
mostró que omitir retornos de exclusión puede distorsionar la medición. En este proyecto el problema
adopta otra forma: la pertenencia histórica se conoce, pero el proveedor gratuito deja de servir
precios de numerosos símbolos retirados. El universo correcto y el panel observable son, por tanto,
dos poblaciones diferentes. Esta distinción justifica que la cobertura se mida contra los miembros
reales del índice y no contra las filas que el propio pipeline consiguió construir.

En síntesis, la contribución del TFM no consiste en proponer un sexto factor. Consiste en evaluar una
arquitectura de aprendizaje con una separación explícita entre datos disponibles, modelo elegido,
robustez y cartera. Ese posicionamiento condiciona todos los capítulos posteriores.

\section{Medición de señal, amplitud y transferencia: Rank-IC como criterio de selección}

La métrica elegida para seleccionar el sistema es la correlación de rangos y no el error cuadrático
de una predicción puntual. Para seleccionar acciones importa más ordenar correctamente dos empresas
que estimar con precisión el retorno porcentual de cada una. La correlación de Spearman es además
robusta a transformaciones monótonas del score: si un modelo cambia la escala pero no el orden, su
Rank-IC no cambia. Un Rank-IC no debe presentarse aislado: su estabilidad se describe mediante la
fracción de cohortes positivas, su dispersión, intervalos por bloques y el IC-IR, definido como la
media de Rank-IC dividida por su desviación típica. Cuando las etiquetas se solapan, el número
nominal de cohortes sobrestima la cantidad de información independiente; el uso de correcciones
Newey--West y bloques de doce meses reconoce precisamente esa dependencia.

La literatura de gestión activa distingue además señal de implementación, y conviene enunciar aquí
esa relación clásica --- la ley fundamental de la gestión activa de Grinold (1989) ---, aunque este
trabajo la use como marco conceptual y no como identidad exacta. El \textit{information ratio}
alcanzable por un gestor depende aproximadamente de tres factores: la calidad de su señal, medida
por el \textit{information coefficient}; la amplitud, es decir, el número de decisiones
aproximadamente independientes que puede tomar; y el coeficiente de transferencia, que mide qué
proporción de esa señal llega efectivamente a la cartera.

\[
  \text{IR} \;\approx\; \text{IC} \times \sqrt{\text{amplitud}} \times \text{TC}
\]

Clarke et al. (2002) muestran por qué las restricciones reducen ese coeficiente de transferencia:
una señal puede contener información que una cartera \textit{long-only}, concentrada o sometida a
límites de rotación no consigue expresar. El punto es central en este TFM. El Rank-IC evalúa toda la
sección transversal, incluida la mitad inferior que una cartera sin posiciones cortas nunca puede
monetizar, mientras que la cartera mantiene ocho nombres y puede retenerlos doce meses. El salto
entre ambas métricas no es una conversión automática, sino un segundo problema de decisión.

Los costes añaden otra ruptura. Amihud (2002) relaciona iliquidez y retornos, y Novy-Marx y Velikov
(2016) muestran que el turnover puede consumir buena parte de las anomalías más intensivas en
negociación. Restar una comisión fija al final no basta cuando el coste también debe cambiar la
decisión de operar. Por eso esta memoria distingue la ruta congelada ---qué habría costado ejecutar
las mismas órdenes--- de la ruta resimulada, en la que un coste mayor eleva los umbrales y puede
evitar la operación.

Su utilidad aquí es que separa tres causas distintas de un mal resultado que una cifra de
rentabilidad confundiría: que la señal no ordene, que se tomen pocas decisiones independientes o que
la cartera no consiga expresar lo que la señal indica. Y tres restricciones de este diseño actúan
directamente sobre el tercer término: la cartera es \textit{long-only}, de modo que la mitad de la
información que el Rank-IC mide es por construcción inaccesible; contiene unas pocas decenas de
posiciones sobre un universo de unos quinientos valores; y paga costes proporcionales a su rotación.

De ahí que la métrica de selección sea el Rank-IC y no el CAGR: seleccionar por rentabilidad final
confundiría los tres términos, mientras que medir la señal aislada permite atribuir después, con un
experimento controlado, cuánto pierde la implementación. El trabajo usa esta descomposición como
marco conceptual y no como identidad: el Capítulo~\ref{chap:resultados-economicos} declara que el
coeficiente de transferencia no está disponible para la cartera adoptada y mide la pérdida de
transmisión por otra vía.

\section{Posición de este trabajo}

La tabla anterior resume cómo responde este trabajo a los riesgos de cada línea. Queda una tensión
que no cabe en una celda y que conviene declarar: la que hay entre señal y utilidad. Buena parte de
la literatura reporta capacidad predictiva sin llevarla a una cartera con costes, o reporta
rentabilidad sin aislar la señal. Aquí se hacen ambas cosas, pero en un orden que importa: primero
la señal, después la cartera, y nunca al revés.

Ese orden permite una decisión metodológica que la literatura suele tratar como contradictoria: la
búsqueda predictiva es secuencial e iterada, y la de cartera cartesiana y optimizada por Information
Ratio, siguiendo la distinción entre calidad de señal e implementación que formula López de Prado
(2018). Por qué ambas cosas son compatibles se argumenta en la
Sección~\ref{sec:seleccion-secuencial}. Lo que ninguna de las dos evita es la multiplicidad frente a
la afirmación económica, y por eso el trabajo reporta el Deflated Sharpe (Bailey y López de Prado,
2014) aun sabiendo que su propia estrategia de búsqueda lo empeora.

La posición final puede resumirse en cuatro decisiones. Primera, el objeto primario es una
ordenación, no una predicción puntual de rentabilidad. Segunda, todo aprendizaje y toda combinación
se realizan con información PIT y etiquetas cerradas. Tercera, el modelo se elige sin observar su
cartera y la cartera se optimiza después sin poder reescribir el modelo. Cuarta, la evidencia
negativa ---un agente que supera al meta, un Deflated Sharpe insuficiente o una reserva débil--- se
conserva con el mismo rango que los contrastes favorables. La contribución no consiste en eliminar
las tensiones de la literatura, sino en construir un protocolo en el que cada una deje una huella
auditable.
